\subsection{取值于 Banach 空间中的函数的积分}

定理 1（Lebesgue--Fubini）. --- 设 f 是定义在 T × T′ 上、取值于一个 Banach 空间 F 或 \(\overline{R}\) 中的函数；设 N 是由这样的 \(t \in T\) 组成之集：函数 \(t' \mapsto \mathbf{f}(t, t')\) 不是 \(\mu'\) -可积。

\begin{enumerate}
\def\labelenumi{\alph{enumi})}
\tightlist
\item
  设 \(\mathbf{f}\) 是 \(\nu\) -可积的；则 \(\mathrm{N}\) 是 \(\mu\) -可忽略的，函数 \(t \mapsto \int \mathbf{f}(t, t') d\mu'(t')\) （对 \(t \notin \mathrm{N}\) 有定义）是 \(\mu\) -可积的，且
\end{enumerate}

\[
\iint \mathbf {f} (t, t ^ {\prime}) d \mu (t) d \mu^ {\prime} (t ^ {\prime}) = \int d \mu (t) \int \mathbf {f} (t, t ^ {\prime}) d \mu^ {\prime} (t ^ {\prime}).\tag{13}
\]

\begin{enumerate}
\def\labelenumi{\alph{enumi})}
\setcounter{enumi}{1}
\tightlist
\item
  设 \(\mathbf{f}\) 是本质 \(\nu\) -可积的，且测度 \(\mu'\) 是适度的；则 \(\mathbf{N}\) 是局部 \(\mu\) -可忽略的，函数 \(t \mapsto \int \mathbf{f}(t, t') d\mu'(t')\) （对 \(t \notin \mathbf{N}\) 有定义）是本质 \(\mu\) -可积的，且 (13) 成立。
\end{enumerate}

断言 a) 由 §3, No.~3 的定理 1 立即得出。为建立 b)，用 g 表示一个局部几乎处处等于 f 的 \(\nu\) -可积函数，并用 H 表示由这样的 \((t, t')\) 组成之集：\(\mathbf{f}(t, t') \neq \mathbf{g}(t, t')\) 。依命题 7 的推论 1，截口 \(\mathrm{H}(t)\) 是 \(\mu'\) -可忽略的，除去 \(t \in T\) （组成一个局部 \(\mu\) -可忽略集者）外。于是把 a) 应用于 g 即得关于 f 的结论。

概说. --- 设 f 是定义在 \(T \times T'\) 上、取值于 \(\overline{R}\) 或一个 Banach 空间中的函数，且是 \(\nu\) -可测和 \(\nu\) -适度的。为使三个积分

\[
\iint \mathbf {f} (t, t ^ {\prime}) d \mu (t) d \mu^ {\prime} (t ^ {\prime}), \int d \mu (t) \int \mathbf {f} (t, t ^ {\prime}) d \mu^ {\prime} (t ^ {\prime}), \int d \mu^ {\prime} (t ^ {\prime}) \int \mathbf {f} (t, t ^ {\prime}) d \mu (t)
\]

都存在并且相等，必须且只须两个数

\[
\int^ {*} d \mu (t) \int^ {*} | {\bf f} (t, t ^ {\prime}) | d \mu^ {\prime} (t ^ {\prime}), \quad \int^ {*} d \mu^ {\prime} (t ^ {\prime}) \int^ {*} | {\bf f} (t, t ^ {\prime}) | d \mu (t)
\]

之一为有限。

这是定理 1、命题 7 和可积性判别法（Ch. IV, §5, No.~6, 定理 5）的直接推论。

注记. --- 1) 当测度 \(\mu'\) 不是适度的时候，可能发生这样的情形：\(\mathbf{f}\) 本质 \(\nu\) -可积，而函数 \(t' \mapsto \mathbf{f}(t, t')\) 不是本质 \(\mu'\) -可积的（对任何 \(t \in \mathrm{T}\) ）（§3, 习题 4）。

\begin{enumerate}
\def\labelenumi{\arabic{enumi})}
\setcounter{enumi}{1}
\tightlist
\item
  设 \(\mu\) 和 \(\mu'\) 是两个复测度，且 \(\nu = \mu \otimes \mu'\) 。若 \(\mathbf{f}\) 是 \(\nu\) -可积的（换句话说，\(|\nu|\) -可积的），则把本定理应用于测度 \(|\mu|\) 和 \(|\mu'|\) （它们的乘积是 \(|\nu|\) ，Ch. III, §4, No.~2, 命题 3）就蕴涵：\(t' \mapsto \mathbf{f}(t, t')\) 是 \(\mu'\) -可积的（对 \(\mu\) -几乎所有的 \(t\) ）。由此，把测度 \(\mu\) 和 \(\mu'\) 分解成正测度的线性组合，即得 a) 的陈述推广到复测度。对 b) 可类似论证。
\end{enumerate}

命题 9. --- 设 F、F′ 和 G 是三个 Banach 空间，并设 \((\mathbf{x},\mathbf{y})\mapsto[\mathbf{x}\cdot\mathbf{y}]\) 是 \(F\times F'\) 到 G 中的一个连续双线性映射。设 f（相应地，\(f'\) ）是定义在 T（相应地，\(T'\) ）上、取值于 F（相应地，\(F'\) ）中且关于 \(\mu\) （相应地，\(\mu'\) ）本质可积的函数。用 g 表示函数 \((t,t')\mapsto[\mathbf{f}(t)\cdot\mathbf{f}'(t')]\) ；则 g 关于 \(\mu\otimes\mu'\) 本质可积，且

\[
\iint \left[ \mathbf {f} (t) \cdot \mathbf {f} ^ {\prime} \left(t ^ {\prime}\right) \right] d \mu (t) d \mu^ {\prime} \left(t ^ {\prime}\right) = \left[ \left(\int \mathbf {f} d \mu (t)\right) \cdot \left(\int \mathbf {f} ^ {\prime} \left(t ^ {\prime}\right) d \mu^ {\prime} \left(t ^ {\prime}\right)\right) \right].\tag{14}
\]

此外，若 \(\mathbf{f}\) 和 \(\mathbf{f}'\) 可积，则 \(\mathbf{g}\) 可积。

函数 \((t, t') \mapsto [\mathbf{f}(t) \cdot \mathbf{f}'(t')]\) 依命题 3 的推论 1 是 \((\mu \otimes \mu')\) -可测的。另一方面，若用 \(b\) 表示双线性映射 \((\mathbf{x}, \mathbf{y}) \mapsto [\mathbf{x} \cdot \mathbf{y}]\) 的范数，则

\[
\begin{array}{r l} \iint^ {\bullet} | [ \mathbf {f} (t) \cdot \mathbf {f} ^ {\prime} (t ^ {\prime}) ] | d \mu (t) d \mu^ {\prime} (t ^ {\prime}) & \leqslant b \iint^ {\bullet} | \mathbf {f} (t) | \cdot | \mathbf {f} ^ {\prime} (t ^ {\prime}) | d \mu (t) d \mu^ {\prime} (t ^ {\prime}) \\ & = b \left(\int^ {\bullet} | \mathbf {f} (t) | d \mu (t)\right) \left(\int^ {\bullet} | \mathbf {f} ^ {\prime} (t ^ {\prime}) | d \mu^ {\prime} (t ^ {\prime})\right) \end{array}
\]

这依据命题 8。这表明 \([\mathbf{f}(t) \cdot \mathbf{f}'(t')]\) 关于 \(\mu \otimes \mu'\) 本质可积（§1, No.~3, 命题 9）。设 \(\mathbf{f}\) 和 \(\mathbf{f}'\) 可积：则 \(\mathbf{f}\) 和 \(\mathbf{f}'\) 是适度的，且 \(\mathbf{g}\) 是适度的（命题 5 的推论 2），从而可积（§1, No.~3, 命题 9 的推论）。在这一情形，公式 (14) 由 Lebesgue--Fubini 定理及积分的线性性得出（Ch. IV, §4, No.~2, 定理 1）。为完成 \(\mathbf{f}\) 和 \(\mathbf{f}'\) 本质可积情形的处理，把 (14) 应用于两个可积函数 \(\mathbf{f}_1\) 和 \(\mathbf{f}_1'\) （分别局部几乎处处等于 \(\mathbf{f}\) 和 \(\mathbf{f}'\) ），并注意 \([\mathbf{f} \cdot \mathbf{f}'] = [\mathbf{f}_1 \cdot \mathbf{f}_1']\) 在 \(T \times T'\) 中局部几乎处处成立（命题 4）。

这一结果推广到复测度的乘积。

